\documentclass[11pt]{article}
\usepackage[margin=1in]{geometry}
\usepackage{amsmath,amssymb,amsthm}
\usepackage[colorlinks=true,urlcolor=blue]{hyperref}
\newtheorem{theorem}{Theorem}
\newtheorem{lemma}{Lemma}
\newtheorem{corollary}{Corollary}
\newcommand{\one}{\mathbf 1}
\newcommand{\Prob}{\mathbb P}
\title{Prescribed versus adaptive cross matchings:\\
a sharp independence boundary for planted RSBM eigenvalue simplicity}
\author{Partitioning-limits session; independent audit supplement}
\date{November 15, 2026}
\begin{document}
\maketitle
\begin{abstract}
Let the two internal graphs be independent uniform simple cubic graphs
on even block size $N$. An adaptive simple perfect matching between
the blocks can force a second balanced $(3,1)$-equitable partition with
probability tending to one. The extra adjacency eigenvector of eigenvalue
two can have support only $2\lfloor\log\log N\rfloor$. All three graph
factors can even be pairwise independent with their usual marginal
distributions. This distinguishes
uniformity over every \emph{prescribed} cross graph from robustness to
cross graphs chosen after observing the internal graphs. The short-cycle
input is classical; no historical priority is asserted for this boundary.
\end{abstract}

Throughout, $C,D$ denote adjacency matrices of labelled simple cubic
graphs on disjoint blocks $V_+,V_-$ of size $N$. A cross matching has
biadjacency matrix $B$, and
\[
 A_B=\begin{pmatrix}C&B\\B^T&D\end{pmatrix},\qquad
 \sigma=(\one_N,-\one_N)^T.
\]
Every such graph has $A_B\sigma=2\sigma$. A balanced $(3,1)$-equitable
partition assigns three same-side neighbors and one opposite-side
neighbor to every vertex, with exactly $N$ vertices on each side.

\begin{lemma}[Clean growing cycles]\label{lem:cycle}
Let $H_N$ be a uniform simple cubic graph on even $N$ labelled vertices.
Put $\ell=\lfloor\log\log N\rfloor$, using natural logarithms. With
probability tending to one, $H_N$ has an induced $\ell$-cycle $X$ such
that its $\ell$ third neighbors outside $X$ are distinct. In particular,
the map sending a cycle vertex to its unique outside neighbor is a
bijection from $X$ to a boundary set $U$ of size $\ell$.
\end{lemma}
\begin{proof}
McKay--Wormald--Wysocka~\cite[Theorem 2]{MWW} applies directly to the
uniform simple labelled model when $(d-1)^{2g-1}=o(N)$. Set $d=3$,
$g=\ell$, and take the set of forbidden lengths to be $\{\ell\}$.
Since $2^{2\ell-1}=o(N)$ and $2^\ell/(2\ell)\to\infty$, its conclusion
gives
\[
 \Prob(H_N\text{ has no }\ell\text{-cycle})
             =\exp\{-2^\ell/(2\ell)+o(1)\}=o(1).
\]
Moreover, Equation (2.9) of the same paper gives probability
$O(2^{2\ell-1}/N)=o(1)$ that an $\ell$-cycle shares an edge with a
different cycle of length at most $\ell$.

For $\ell\ge4$, an $\ell$-cycle with a chord shares an arc with a
different cycle of length at most $\lfloor\ell/2\rfloor+1\le\ell$.
If two of its vertices have the same outside third neighbor, the
shorter arc between them, together with the two boundary edges,
forms a different cycle of length at most
$\lfloor\ell/2\rfloor+2\le\ell$, again sharing a cycle edge.
Thus, outside an event of probability $o(1)$, every $\ell$-cycle is
clean in the stated sense and at least one exists.
\end{proof}

\begin{lemma}[Deterministic adaptive matching construction]
\label{lem:construction}
Suppose $C$ has a clean $\ell$-cycle $X\subset V_+$ with outside
boundary $U$, and $D$ has a clean $\ell$-cycle $Y\subset V_-$ with
outside boundary $W$, with $2\ell\le N$ and $0<\ell<N$.
Choose any bijections matching $X$ to $W$, $U$ to $Y$, and
$V_+\setminus(X\cup U)$ to $V_-\setminus(Y\cup W)$.
Their union is a simple perfect matching $B$. Then
\[
 \zeta=\binom{\one_X}{-\one_Y}\quad\text{satisfies}\quad
 A_B\zeta=2\zeta.
\]
The vectors $\sigma,\zeta$ are linearly independent. Furthermore
$\sigma'=\sigma-2\zeta$ is the sign vector of a second balanced
$(3,1)$-equitable partition, distinct from $\pm\sigma$.
\end{lemma}
\begin{proof}
On $X$, the two cycle neighbors contribute two to $A_B\zeta$, and
the outside neighbor and cross neighbor lie in $U,W$ where $\zeta=0$.
On $U$, its unique neighbor in $X$ contributes one while its matching
neighbor in $Y$ contributes minus one; its other internal neighbors
are outside $X$, so the total is zero. The corresponding calculation
on $Y,W$ gives minus two and zero. Vertices outside these four sets
have no internal neighbor in $X$ or $Y$, and their cross neighbors
also lie outside $X,Y$, so their total is zero. This proves the
eigenvector identity.

Since $\zeta$ vanishes at vertices outside $X\cup Y$ and is nonzero,
it cannot be a scalar multiple of the everywhere nonzero $\sigma$.
The vector $\sigma'$ flips precisely $X$ and $Y$. These sets have
equal size, so the new partition is balanced. It is a sign vector,
and linearity gives $A_B\sigma'=2\sigma'$. Because every vertex
has total degree four, this equation says its same-side degree is
$(4+2)/2=3$ and its opposite-side degree is $(4-2)/2=1$.
Finally $0<\ell<N$ rules out both global sign choices of $\sigma$.
\end{proof}

\begin{theorem}[Adaptive failure with high probability]
\label{thm:adaptive}
Let $C,D$ be independent uniform simple cubic graphs on even $N$.
There is a deterministic function $F_N(C,D)$ taking values in the
simple perfect matchings between their blocks such that
\[
 \Prob\!\left(\dim\ker_{\mathbb R}
      (A_{F_N(C,D)}-2I)\ge2\right)\longrightarrow1.
\]
On an event of probability tending to one the graph has two distinct
balanced $(3,1)$-equitable partitions, and one additional eigenvector
has support $2\lfloor\log\log N\rfloor$.
The function can be evaluated in polynomial time.
\end{theorem}
\begin{proof}
Choose the first clean $\ell$-cycle in each graph in a fixed ordering,
and apply Lemma~\ref{lem:construction}, resolving its bijections in a
fixed ordering. If either graph has no such cycle, choose an arbitrary
fixed cross matching. Lemma~\ref{lem:cycle} and a union bound imply
that both cycles exist with probability tending to one. The deterministic
lemma proves the stated conclusions.

In a cubic graph the number of non-backtracking walks of length
$\ell$ starting from all vertices is at most $3N2^{\ell-1}$.
Enumerating these walks, checking closed simple cycles and their
boundaries, takes $O(N\ell2^\ell)$ operations in a standard adjacency-list
model. For $\ell=\lfloor\log\log N\rfloor$ this is polynomial in $N$.
The remaining matching construction also takes polynomial time.
\end{proof}

\begin{corollary}[Pairwise independence is insufficient]
\label{cor:marginal}
There is a joint law of $(C,D,B)$ such that the three factors are
pairwise independent, $C,D$ each are uniform simple cubic graphs,
and $B$ is a uniform simple perfect matching, but
\[
             \Prob(\dim\ker(A_B-2I)\ge2)\longrightarrow1.
\]
Thus even the same pairwise laws as the independent RSBM do not
imply planted-eigenvalue simplicity. Independence of $B$ from the
joint pair $(C,D)$ is essential to the positive theorem.
\end{corollary}
\begin{proof}
View a cross matching as a bijection $m:V_+\to V_-$. Independently
sample uniform permutations $\pi,\rho$ of the two blocks, relabel the
internal graphs by $C'=\pi C\pi^{-1}$, $D'=\rho D\rho^{-1}$, and let
$m'=F_N(C',D')$. Return the matching
\[
                              m=\rho^{-1}m'\pi
\]
on the original labels. The internal graph law is invariant under
all label permutations. Consequently $(C',D')$ is independent of
$(\pi,\rho)$, so $m'$ is independent of them. Conditional on any
$m'$, the map $\rho^{-1}m'\pi$ is a uniform bijection. Hence $B$
has the claimed marginal.

For the stronger independence assertion, condition on $C,\pi,D'$.
The value $m'=F_N(\pi C\pi^{-1},D')$ is then fixed, whereas $\rho$
is still independent and uniform: the relabelled uniform graph $D'$
is independent of $\rho$. Hence $m=\rho^{-1}m'\pi$ is uniform
conditional on $C$, proving $B$ independent of $C$. Conditioning on
$D,\rho,C'$ proves $B$ independent of $D$ in the same way. The pair
$(C,D)$ was never changed and retains its independent uniform law.
Relabelling preserves the extra eigenvector and partition, so
Theorem~\ref{thm:adaptive} proves the failure claim.
\end{proof}

\paragraph{The quantifier boundary.}
The positive prescribed-$B$ theorem has the form
\[
       \sup_{B\text{ fixed}}\Prob_{C,D}(\text{failure for }B)=o(1).
\]
The result here is
\[
       \Prob_{C,D}(\exists B:\text{failure for }B)=1-o(1).
\]
These are compatible. The latter refutes an adaptive strengthening,
and Corollary~\ref{cor:marginal} shows that even pairwise independence
does not supply the required joint independence. The random construction changes only
$2\ell=o(N)$ labels, so no macroscopic disagreement claim is made.
This proves failure of eigenvalue simplicity and combinatorial partition
uniqueness; it does not by itself prove a statistical recovery lower
bound for an estimator knowing the adaptive generative law.
The use of a classical short-cycle theorem is explicit; this note is
not a claim of historical firstness or of a new random-graph method.

\begin{thebibliography}{9}
\bibitem{MWW}
B.~D.~McKay, N.~C.~Wormald and B.~Wysocka,
\emph{Short cycles in random regular graphs},
Electronic Journal of Combinatorics \textbf{11} (2004), Research Paper 66.
Theorem 2 and Equation (2.9).
\url{https://users.monash.edu/~nwormald/papers/girth.pdf}.
\end{thebibliography}
\end{document}
